%! Introducing Statistics: What Can We Learn from Data? — Unit 1, Lesson 1.0 — Experience & Formalize
%! (Activity · QuickNotes · Application · Check Your Understanding)
%! Directory stays `experience/` (a build identifier in shared/lesson.mk); the label
%! students and teachers read is "Experience & Formalize".
% EFFL component, Math Medic sizing: 12pt body, OPEN white space after every question.
% \answerspace{H}{} reserves exactly H in the blank; the key passes the answer as the second
% argument, so the two files paginate identically by construction. Keep the macro
% byte-identical in the blank and the key. \boxguard counts here are 12pt-relative (~14-16).
% Page budget: Activity <=2pp · QuickNotes 1/2pp · Application 1/2-1pp · CYU 1-2pp.
%
% SPOILER RULE: the Activity never names statistics / data / individual / variable / value /
% variability / statistical question. Students answer in their own words ("row", "column
% heading", "look it up", "read the whole column"); the debrief attaches the formal terms in
% QuickNotes. Nothing above the QuickNotes box may use the vocabulary.
%
% This lesson has no computation, so it carries no \begin{work} block — page lockstep here
% rests entirely on \answerspace. (CED Topic 1.1 / VAR-1.A: identifying questions, not computing.)
\documentclass[12pt]{article}
\usepackage{apstats-article}
\usepackage{apstats-boxes}
\newcommand{\answerspace}[2]{\par\nopagebreak\noindent\begin{minipage}[t][#1][t]{\linewidth}%
  \color{keyred}\bfseries #2\end{minipage}\par}

\begin{document}

\pageheader{Unit 1, Lesson 1.0}{Experience \& Formalize: The Adoption Board}
% NAMESTRIP: no \namedateperiod here — the name row lives on the cover only.

\vspace{0.06in}

% ── Part 1: the Activity ─────────────────────────────────────────────────────
% Scenario 1 builds the toolkit on one pre-drawn table. Scenario 2 is the contrast case and
% carries the CRUX: "How many dogs are on the board?" needs the whole table AND has a number,
% but is NOT statistical, because the answer does not differ from one dog to the next. That
% single discriminator — do the answers vary across individuals? — is the lesson.
\begin{headlinebox}{sky}
\textbf{The Adoption Board.} A county animal shelter posts one board in its lobby listing every
dog currently up for adoption. Visitors read it, and so does the shelter director --- but they
are not looking for the same things. With your group, work through the two problems below using
only what you already know, and be ready to explain \emph{how you know} each answer
\textbf{from the board itself}.
\end{headlinebox}

\vspace{0.04in}

\begin{scenariobox}[1. Reading the Board]{navy}
\small
\renewcommand{\arraystretch}{1.25}
\begin{tabularx}{\linewidth}{@{} l Y Y Y Y @{}}
\rowcolor{sky}
\textbf{Name} & \textbf{Breed} & \textbf{Age (yr)} & \textbf{Weight (lb)} & \textbf{Days on board} \\
\hline
Biscuit & Beagle    & 3 & 24 & 12 \\
Momo    & Terrier   & 1 & 15 & 44 \\
Ranger  & Shepherd  & 5 & 61 &  6 \\
Pepper  & Retriever & 2 & 55 & 21 \\
Sage    & Beagle    & 7 & 27 & 59 \\
Tuck    & Mixed     & 4 & 38 &  9 \\
Willow  & Terrier   & 2 & 17 & 33 \\
Zeke    & Shepherd  & 6 & 70 & 15 \\
\hline
\end{tabularx}

\normalsize
\begin{enumerate}[label=\textbf{\alph*.}, leftmargin=1.8em, itemsep=8pt]
  \item One \textbf{row} of this board describes \blank{5.2cm}.

        The board has \blank{1.4cm} rows and \blank{1.4cm} columns of recorded facts
        (not counting the Name column).
  \item ``Beagle'' appears on the board. Is it a \textbf{column heading}, or something
        \textbf{written underneath} one? Circle your answer, then say which column heading it
        belongs to. \hfill \textbf{heading} \qquad \textbf{written underneath}
        \answerspace{1.4cm}{}
  \item \textbf{Underline} the part of the board you would read to answer each question.
        One needs a single spot; the other needs a whole column.

        \emph{How much does Zeke weigh?} \qquad \emph{How much do these dogs weigh?}
        \answerspace{1.6cm}{}
  \item Sage has been on the board for 59 days and Ranger for 6. If \emph{every} dog had been
        on the board the same number of days, would the ``Days'' column still be worth posting?
        Explain.
        \answerspace{2.0cm}{}
\end{enumerate}
\end{scenariobox}

\boxguard[16]
\begin{scenariobox}[2. Four Questions the Director Might Ask]{navy}
The shelter director looks at the same board and asks these four questions.

\begin{enumerate}[label=\textbf{\alph*.}, leftmargin=1.8em, itemsep=8pt]
  \item For each one, decide whether you could answer it by reading \textbf{one spot} on the
        board, or whether you would have to read \textbf{many rows}. Write \emph{one spot} or
        \emph{many rows} in the blank.

        \begin{enumerate}[label=\textbf{(\Alph*)}, leftmargin=2.2em, itemsep=5pt]
          \item Is Momo a terrier? \hfill \blank{3.2cm}
          \item How many dogs are on the board? \hfill \blank{3.2cm}
          \item How long do dogs on this board tend to wait? \hfill \blank{3.2cm}
          \item Do heavier dogs tend to wait longer than lighter ones? \hfill \blank{3.2cm}
        \end{enumerate}
  \item Two of those needed many rows. But only \emph{one} of the two has an answer that
        \textbf{changes from one dog to the next}. Which one is which, and how can you tell?
        \answerspace{2.6cm}{}
  \item A flyer taped under the board reads: \textbf{``Our dogs go home in 25 days.''}
        Name \textbf{two} things the flyer does not tell you that you would want to know
        before believing it.
        \answerspace{2.0cm}{}
  \item A visitor reads the flyer and says, ``Good --- so I'll have my dog home in 25 days.''
        Using the board, give one specific reason that is the wrong way to read the number.
        \answerspace{2.0cm}{}
\end{enumerate}
\end{scenariobox}

% ── Part 2: QuickNotes (the debrief fills this) — target: half a page ────────
\boxguard[14]
\begin{tcolorbox}[enhanced, colback=sky, colframe=navy, boxrule=0.8pt, arc=2mm,
    left=3mm, right=3mm, top=1.5mm, bottom=1.5mm, breakable,
    fonttitle=\bfseries\small\color{white}, title={QuickNotes: Data, Variables, and Statistical Questions},
    attach boxed title to top left={yshift=-2mm, xshift=4mm},
    boxed title style={colback=navy, colframe=navy, arc=1.5mm}]
\small
\begin{minipage}[t]{0.36\linewidth}
\vspace{2pt}
\renewcommand{\arraystretch}{1.15}
\begin{tabular}{@{}l l@{}}
\multicolumn{2}{@{}l}{\footnotesize\textbf{\color{navy}Breed} \quad \textbf{\color{navy}Weight}}\\
\hline
Beagle    & 24 \\
Terrier   & 15 \\
Shepherd  & 61 \\
\hline
\end{tabular}\\[3pt]
{\footnotesize Headings are \textbf{variables}.\\
What is written under them\\
are \textbf{values}. Each row is\\
one \textbf{individual}.}
\end{minipage}\hfill
\begin{minipage}[t]{0.60\linewidth}
\vspace{-4pt}
\begin{itemize}[leftmargin=1.1em, itemsep=5pt]
  \item \textbf{Data} are recorded values \emph{plus} the \blank{2.4cm} that gives them meaning.
        A number alone answers nothing.
  \item Before a number is data, answer the \textbf{W questions}: who, \blank{1.6cm},
        in what \blank{1.8cm}, when and where.
  \item An \textbf{individual} is what one \blank{1.6cm} describes. A \textbf{variable} is a
        \blank{2.0cm} heading --- never one of the entries under it.
  \item \textbf{Variability:} values \blank{2.4cm} from one individual to the next.
        \emph{This is why statistics exists.}
  \item A \textbf{statistical question} passes \textbf{both} checks:\\[2pt]
        \textbf{1.} Do the answers \blank{2.2cm} across individuals?\\[2pt]
        \textbf{2.} Must I \blank{2.6cm} data to answer it?
\end{itemize}
\end{minipage}
\end{tcolorbox}

% ── Part 3: the Application (worked together, ~6 min) ────────────────────────
% No \begin{work} block: this topic has no computation. Page lockstep rests on \answerspace.
\boxguard[14]
\begin{notesbox}{Application: The Gym Sign-Up}
We will do this one together. A gym records the \textbf{resting heart rate}, in beats per
minute, of each of its 120 members on the morning they join. It also records each member's
\textbf{home ZIP code}.

\begin{enumerate}[label=\textbf{\alph*.}, leftmargin=1.8em, itemsep=7pt]
  \item The individuals are \blank{6.4cm}, and two variables recorded are
        \blank{4.0cm} and \blank{4.0cm}.
  \item Write one statistical question this gym could answer, and say exactly what varies in it.
        \answerspace{1.8cm}{}
  \item \textbf{What if we change the variable?} ``What is the most common ZIP code among the
        members?'' --- is that a statistical question? Run both checks before you answer.
        \answerspace{2.2cm}{}
\end{enumerate}
\end{notesbox}

% ── Part 4: Check Your Understanding — practice, NOT scored ──────────────────
\boxguard[14]
\begin{notesbox}{Check Your Understanding}
Work with your partner, then finish on your own. \emph{These are for practice --- they are not
scored.}
\begin{enumerate}[leftmargin=1.9em, itemsep=9pt]
  \item A biologist weighs each of the 30 turtles in a pond and records the weight in kilograms.

        The individuals are \blank{5.4cm}; the variable is \blank{3.6cm},
        measured in \blank{3.0cm}.
  \item Write one statistical question about those turtles, then name what varies in it.
        \answerspace{1.8cm}{}
  \item \textbf{Same pond, two questions.} Explain why exactly one of these is statistical.

        \emph{(i)} How much does the heaviest turtle weigh? \quad
        \emph{(ii)} How do the turtles' weights vary?
        \answerspace{2.0cm}{}
  \item A report states only: \textbf{``The average is 12.''} Name \textbf{two} pieces of
        context you would need before this number means anything.
        \answerspace{1.8cm}{}
  \item \textbf{Formative check.} Circle the statistical question, then justify it in one
        sentence by pointing to what varies.
        \begin{enumerate}[label=\textbf{(\Alph*)}, leftmargin=2.2em, itemsep=4pt]
          \item How many students are enrolled at this school?
          \item How many hours did Jordan work last week?
          \item How do the hours worked each week vary among students who have jobs?
          \item What time does the library close on Friday?
        \end{enumerate}
        \answerspace{1.6cm}{}
\end{enumerate}
\end{notesbox}

\end{document}
